\section{Prevalence in the literature}

The preceding sections establish, in controlled settings, how a
variance-contribution measure can be misread as an effect size. To estimate how
often that misreading reaches published conclusions, we screened the literature
that uses the \code{ranger} R package --- one of the most widely used
random-forest implementations, and the one whose importances these results
concern. We screen for papers whose stated conclusions depend on the magnitude or
ranking of an impurity importance without the corroboration that would guard
against the effects of Sections~\ref{sec:linear}--\ref{sec:confound}.

\subsection{Method}
\label{sec:prevalence-method}

\paragraph{Corpus.} We took every work recorded by OpenAlex as citing the ranger
reference publication \citep{wright2017} (OpenAlex work \code{W2157395790}),
$3{,}189$ works in total. For each we retrieved bibliographic metadata from
OpenAlex and, where an open-access full text was available, the full text from
Europe PMC; $773$ works had retrievable open-access full text. Screening was
restricted to these, because the detail needed to judge exposure --- which
importance measure is used, and whether it is corroborated --- is almost never
present in an abstract.

\paragraph{Relevance filter.} Many papers cite ranger only for its speed or for
prediction and never interpret a variable importance. To drop these without
discarding relevant work, we required a textual mention of variable importance: a
paper was retained if its full text matched, case-insensitively, any of the
following regular-expression alternatives.
\begin{verbatim}
(variable|feature|permutation|gini|impurity)\s+importance
importance\s+(score|value|ranking|measure)
mean decrease
decrease in (impurity|gini|accuracy)
\bMDI\b   \bMDA\b   \bBoruta\b
impurity[_\s]?corrected
most important (variable|feature|predictor)
\end{verbatim}
$524$ of the $773$ matched. The filter is deliberately broad: an earlier version
that also required a significance keyword (\code{importance\_pvalues}, Altmann,
Janitza, Boruta) missed most exposed papers. Of $40$ sampled papers it rejected,
$24$ still used an impurity importance and $11$ met the flag criterion below,
because many papers rank features by impurity importance with no significance
test. In total $544$ papers were screened: the $524$ matches plus $20$ from
earlier passes (including that probe) that the broadened filter misses. Its recall
was not separately re-estimated (see limitations).

\paragraph{Two-stage screen.} Each paper was read by a large language model under
a fixed instruction and asked to return a structured judgement; the verbatim
prompts are in the supporting repository (Section~\ref{sec:availability}). The
screen has two stages using different models, so that the second acts as an
adversarial check on the first.

\emph{Stage 1 --- screening.} A first model (Claude Haiku 4.5) read each paper and
recorded which importance measure it uses, whether it interprets that importance's
magnitude or ranking, which corroborating analyses it reports, and a flag for
whether a central conclusion could have been affected. The instruction was
explicit that the bias corrupts magnitude and ranking, not the binary verdict
that a variable is used, so a significant importance $p$-value
\citep{altmann2010,janitza2018,kursa2010}
does not by itself protect a conclusion. A paper was flagged only if it leaned on
the magnitude or ranking of an impurity importance for a central conclusion
\emph{without} corroboration, defined as permutation importance,
SHAP \citep{lundberg2017}, partial-dependence/ALE \citep{friedman2001,apley2020},
conditional importance \citep{strobl2008}, or held-out validation of the
\emph{selected} features (overall test accuracy does not count, as it validates
prediction, not the ranking). Every flag required a quoted passage.

Counting SHAP as corroboration is generous. Section~\ref{sec:confound} shows
that a global mean $|\text{SHAP}|$ ranking carries the same shape/skew confound
as impurity importance, while a SHAP \emph{dependence} plot does not, and the
screen does not distinguish the two forms, so papers reporting only a mean
$|\text{SHAP}|$ ranking were credited with corroboration they do not have. This
can only suppress flags, so it makes the confirmed count conservative.

\emph{Stage 2 --- adjudication.} Every flagged paper was independently re-read by
a second, stronger model (Claude Sonnet 5) under an explicitly adversarial
instruction: its task was to \emph{refute} the flag. It returned a verdict of
\emph{confirmed}, \emph{refuted}, or \emph{uncertain}, and confirmed only when all
three conditions held at once: 1) the paper genuinely uses an impurity (not solely
a permutation or SHAP) importance, 2) the ranking is central to a stated conclusion,
and 3) the ranking is not adequately corroborated. Both models were accessed in
July 2026. Using two models from a single vendor means the adjudication is a
precision filter, not an independent validation --- a different-vendor or human
check would be stronger; we return to this in the limitations.

\subsection{Results}
\label{sec:prevalence-results}

Table~\ref{tab:prevalence} traces the corpus through the pipeline. Of the $544$
papers screened, $307$ were judged to use an impurity importance and $125$ were
flagged in stage~1. The adversarial adjudication confirmed $55$, refuted $65$, and
could not decide for $5$, overturning roughly $56\%$ of the stage-1 flags. The
implied stage-1 precision is therefore about $44\%$, and raw screening counts
should be discounted accordingly.

\begin{table}[htb]
\centering
\caption{The ranger-citing corpus traced through the screening pipeline. The
final rows are counts of screened papers; percentages in the text are relative to
the $544$ screened.}
\label{tab:prevalence}
\begin{tabular}{lr}
\toprule
Stage & Count \\
\midrule
Works citing ranger (OpenAlex) & $3{,}189$ \\
With open-access full text & $773$ \\
Matching the relevance filter & $524$ \\
Screened (filter matches $+$ non-match sample) & $544$ \\
\midrule
Using an impurity importance (stage 1) & $307$ \\
Flagged (stage 1, Haiku 4.5) & $125$ \\
\quad Confirmed (stage 2, Sonnet 5) & $55$ \\
\quad Refuted & $65$ \\
\quad Uncertain & $5$ \\
\bottomrule
\end{tabular}

\vspace{4pt}
{\footnotesize\emph{Note:} the $544$ screened papers comprise all $524$
filter matches plus $20$ non-matching papers that were screened in earlier passes
before the filter was finalised, including the random recall sample; all $524$
filter matches were screened.}
\end{table}

Taking the confirmed count in Table~\ref{tab:prevalence} at face value, $55$ of
the $544$ screened papers state a central conclusion that could have been
distorted by the effects of Sections~\ref{sec:linear}--\ref{sec:confound} and did
not cross-check it. That is about $10\%$ of the screened papers. Measured against
the $307$ that actually interpret an impurity importance, it is $18\%$, roughly
one in six. As a share of all $773$ full-text ranger-citing works it is about
$7\%$.

Two features of the confirmed set are worth drawing out. First, exposure is
concentrated in papers that use no significance machinery: $51$ of the $55$
confirmed papers ($93\%$) report no importance $p$-value or Boruta selection and
simply rank features by impurity importance, suggesting that familiarity with
significance testing goes with more cautious interpretation of a ranking. Second,
the confirmed papers are not confined to low-visibility venues: citation counts
span the full range of the corpus (median $11$, versus a corpus median of $10$),
and $12$ of the $55$ exceed $50$ citations.

\paragraph{Limitations.} These are model judgements, not audits: a flag means a
conclusion \emph{could} have been affected and was uncorroborated, not that it is
wrong, and predictive accuracy is unaffected. The $10\%$ is an unvalidated
estimate; stage 2 raises precision but is not independent validation, and there
is no human ground truth. Only $773$ of $3{,}189$ works had open-access full
text, and open access correlates with field and venue, so the sample may not
generalise. The broadened filter's recall was not re-estimated. Our controlled
experiments are all regressions, and the screen does not record which of the
screened papers are classification and which regression; the variance reading of
the Gini index in Section~\ref{sec:background-strands} is what carries the
caution across to the classification papers, and we did not verify it there
empirically. The repository
lists flagged papers by identifier; we treat this as a claim about corroboration
practice, not correctness, with each flag backed by a quoted passage.
